% Generated draft: requires mathematical and source review.
\documentclass[11pt]{article}
\usepackage[margin=25mm]{geometry}
\usepackage{fontspec}
\setmainfont{Latin Modern Roman}
\usepackage{amsmath,amssymb,mathtools,unicode-math}
\setmathfont{Latin Modern Math}
\usepackage{enumitem,needspace,xurl,hyperref}
\setlength{\parindent}{0pt}
\setlength{\parskip}{5pt}
\setlength{\emergencystretch}{3em}
\newcommand{\N}{\mathbb{N}}
\newcommand{\Z}{\mathbb{Z}}
\newcommand{\Q}{\mathbb{Q}}
\newcommand{\R}{\mathbb{R}}
\newcommand{\C}{\mathbb{C}}
\newcommand{\F}{\mathbb{F}}
\newcommand{\A}{\mathbb{A}}
\newcommand{\PP}{\mathbb{P}}
\newcommand{\E}{\mathbb{E}}
\newcommand{\eps}{\varepsilon}
\newcommand{\dd}{\,\mathrm{d}}
\begin{document}
\section*{UnsolvedMath ID: 30004302}
\textbf{Source problem code:} catalogue entry
\subsection{Definitions and mathematical statement}
Let \(X\) be a set and let \(r:X\times X\to X\times X\) satisfy the set-theoretic Yang--Baxter equation
\[
(r\times\operatorname{id}_X)(\operatorname{id}_X\times r)(r\times\operatorname{id}_X)
=
(\operatorname{id}_X\times r)(r\times\operatorname{id}_X)(\operatorname{id}_X\times r)
\quad\text{on }X^3.
\]
The source does not specify whether its convention for a solution additionally requires bijectivity or nondegeneracy of \(r\), and it does not impose finiteness of \(X\). These qualifications remain unresolved.

Define the structure monoid by
\[
M(X,r)=\left\langle X\ \middle|\ xy=uv\text{ whenever }r(x,y)=(u,v)\right\rangle_{\mathrm{monoid}}.
\]
For a unital coefficient ring \(k\), its monoid algebra \(A=kM(X,r)\) consists of finite formal sums \(\sum_{m\in M(X,r)}a_m m\), with coefficientwise addition and multiplication determined by \((a m)(b n)=(ab)(mn)\). The source does not specify the allowed coefficient rings, whether \(k\) must be a field, or how the sought characterizations may depend on \(k\).

The solution is \emph{involutive} if
\[
r\circ r=\operatorname{id}_{X\times X}.
\]
The monoid \(M=M(X,r)\) is \emph{cancellative} if both cancellation laws hold:
\[
\forall a,b,c\in M,\qquad
ab=ac\Longrightarrow b=c,
\qquad
ba=ca\Longrightarrow b=c.
\]
\textbf{Problem.} Find necessary and sufficient ring-theoretic or homological conditions on the structure algebra \(kM(X,r)\) that characterize, separately, involutivity of \((X,r)\) and cancellativity of \(M(X,r)\). The intended class of solutions and coefficient rings must be specified before these requests become fully quantified characterization problems. No finiteness or nondegeneracy hypothesis is added here to resolve that omission.
\subsection{Short English statement}
Characterize involutivity of a set-theoretic Yang--Baxter solution and two-sided cancellation in its structure monoid through ring-theoretic or homological properties of its structure algebra. The source leaves the coefficient domain and some solution conventions unspecified. Its supplied status is partially solved; this draft does not independently verify that assessment.
\subsection{Sources}
\begin{itemize}
\item Supplied local UnsolvedMath registry record, permanent ID \(30004302\), title: Involutivity and Cancellativity of Yang--Baxter Solutions. Exact source problem code: \texttt{catalogue entry}. Collection: Oberwolfach Reports Open Problems. No individual proposer is named in the supplied record.
\item Supplied original-report citation: \emph{Mini-Workshop: Algebraic Tools for Solving the Yang--Baxter Equation} (2020), Oberwolfach Reports, pp.~3207--3242, DOI \texttt{10.4171/owr/2019/51}. \url{https://doi.org/10.4171/owr/2019/51}. The record gives proposed year 2019; the relationship between that year and the citation's 2020 date is not explained.
\item Supplied literature source: Colazzo--Jespers--Kubat--Van Antwerpen, structure algebras. \url{https://arxiv.org/abs/2305.06023}.
\item Supplied literature source, labeled ``2026 structure-algebra results'': \url{https://doi.org/10.1016/j.jalgebra.2026.04.044}. No further bibliographic details are supplied.
\item Literature assessment checked at \texttt{2026-08-22}: supplied dated curation, not independently verified. It labels the problem partially solved and reports that recent results for finite left-nondegenerate solutions relate involutivity, cancellativity, primeness, domain properties, and homological dimensions, while the requested characterization for arbitrary solutions remains incomplete. The background identifies this as dated literature triage derived from the supplied Oberwolfach Reports curation and requests independent verification of time-sensitive status and cited claims.
\item The registry labels the statement \texttt{corrected\_verified} and reports a manual semantic, typographic, and LaTeX review. These are supplied editorial claims, not independent mathematical verification of this draft or of the literature assessment.
\end{itemize}
\end{document}
